\section{Bottom View of a Binary Tree}
\label{sec:trees-bottom-view}

\problemstatement{The bottom view is the set of nodes visible when the tree
is viewed from directly below: for each horizontal distance (HD), the
\emph{lowest} node on that vertical line (ties broken by the later node in
level order). Return these ordered left to right by HD.}

\begin{patternbox}
Bottom view is top view with one word changed: keep the \emph{last} node
per HD instead of the first. Under BFS -- which scans top-down -- the last
node written to a given HD is the deepest, exactly the one visible from
below. So the code is identical except we \textbf{always overwrite} the
map entry rather than writing only once.
\end{patternbox}

\subsection*{Understanding the problem carefully}

From below, on each vertical line the visible node is the lowest one.
Because BFS visits nodes in non-decreasing depth, repeatedly overwriting
the value stored for an HD leaves, at the end, the value of the last (hence
deepest) node BFS placed there. The left-to-right order within a level
also gives the standard tie-break for two nodes at the same HD and depth.

\begin{ideabox}[BFS, Last Node per HD Wins]
Run the same level-order BFS with HD tags as top view, but write
$\textit{map}[HD]=node.val$ \emph{unconditionally} every time. Later
(deeper) nodes overwrite earlier ones, so each HD ends holding its
bottommost node. Emit ascending by HD.
\end{ideabox}

\subsection*{Algorithm}

\begin{enumerate}
  \item BFS from root with $(node, HD)$; root's HD $=0$.
  \item For each dequeued $(node, HD)$: set $\textit{map}[HD]=node.val$
        (overwrite). Enqueue $(left, HD-1)$, $(right, HD+1)$.
  \item Output the map's values in ascending HD.
\end{enumerate}

\subsection*{Dry run}

\dryrun{Tree: root $1$; left $2$ (right child $4$); right $3$.}

\begin{itemize}
  \item $(1,0)$: map$\{0\!:\!1\}$; enqueue $(2,-1),(3,1)$.
  \item $(2,-1)$: map$\{-1\!:\!2\}$; enqueue $(4,0)$.
  \item $(3,1)$: map adds $1\!:\!3$.
  \item $(4,0)$: overwrite HD $0$ with $4$.
  \item Ascending HD: $2,4,3$ (note HD $0$ now shows $4$, not $1$).
\end{itemize}

\subsection*{C++ implementation}

\begin{lstlisting}
#include <bits/stdc++.h>
using namespace std;

struct TreeNode {
    int val;
    TreeNode* left;
    TreeNode* right;
    TreeNode(int v) : val(v), left(nullptr), right(nullptr) {}
};

vector<int> bottomView(TreeNode* root) {
    vector<int> ans;
    if (root == nullptr) return ans;

    map<int, int> lowAtHD;               // hd -> last (lowest) value
    queue<pair<TreeNode*, int>> q;
    q.push({root, 0});
    while (!q.empty()) {
        auto [node, hd] = q.front();
        q.pop();
        lowAtHD[hd] = node->val;         // always overwrite
        if (node->left)  q.push({node->left,  hd - 1});
        if (node->right) q.push({node->right, hd + 1});
    }
    for (auto& [hd, v] : lowAtHD) ans.push_back(v);
    return ans;
}

int main() {
    TreeNode* root = new TreeNode(1);
    root->left = new TreeNode(2);
    root->right = new TreeNode(3);
    root->left->right = new TreeNode(4);

    for (int x : bottomView(root)) cout << x << " ";
    cout << "\n";
    return 0;
}
\end{lstlisting}

\begin{complexitybox}
\textbf{Time Complexity.} $O(n\log n)$ with an ordered map (or $O(n)$ with
a hash map and a final sort). \textbf{Space Complexity.} $O(n)$ for the
queue and map.
\end{complexitybox}

\begin{takeawaybox}
Top view and bottom view differ by exactly one decision -- keep the first
versus the last node per HD. Recognising that a whole family of ``view''
problems shares one BFS-with-HD engine, differing only in the per-line
selection rule, is the real lesson.
\end{takeawaybox}
